% ============================================================================
% BOLUM 3: BENZER CALISMALAR (LITERATUR TARAMASI)
% ============================================================================
% Kurumsal Temizlik Takip ve Yonetim Sistemi (KTTYS)
% Muhammet Anil YAGIZ - 213255046
% Kirikkale Universitesi - Bilgisayar Muhendisligi Bolumu
% ============================================================================

\chapter{BENZER CALISMALAR}
\label{chap:benzer-calismalar}

\section{Giris}
\label{sec:literatur-giris}

Bu bolumde, kurumsal temizlik yonetimi, is gucu planlamasi ve tesis 
yonetimi alanlarinda yapilmis akademik calismalar ile ticari yazilim 
cozumleri incelenmektedir.

\section{Akademik Calismalar}
\label{sec:akademik-calismalar}

Asagida, proje konusuyla dogrudan veya dolayli olarak iliskili olan 
akademik calismalar incelenmektedir.

\subsection{Akilli Bina Yonetim Sistemleri}
\label{subsec:akilli-bina}

Chen ve Zhang (2019), IoT-Based Smart Building Management System for 
Energy Efficiency and Facility Maintenance baslikli calismalarinda, 
Nesnelerin Interneti (IoT) teknolojilerini kullanarak akilli bina 
yonetim sistemleri gelistirmislerdir.

Arastirmacilar, sistemin uc ana modulden olustugunu belirtmislerdir:
\begin{enumerate}
    \item \textbf{Veri Toplama Modulu:} IoT sensorleri araciligiyla 
    sicaklik, nem, hareket ve enerji tuketimi verilerinin toplanmasi.
    \item \textbf{Analiz Modulu:} Makine ogrenmesi algoritmalari 
    kullanilarak verilerin analiz edilmesi.
    \item \textbf{Yonetim Modulu:} Web tabanli bir arayuz uzerinden 
    bina yoneticilerinin sistemi izlemesi ve yonetmesi.
\end{enumerate}

\subsection{Is Gucu Cizelgeleme Problemleri}
\label{subsec:is-gucu-cizelgeleme}

Ernst ve arkadaslari (2004), Staff scheduling and rostering: A review 
of applications, methods and models baslikli kapsamli derleme 
calismalarinda, personel cizelgeleme problemlerinin matematiksel 
modellemesi ve cozum yontemlerini incelemislerdir.

Calismada, personel cizelgeleme probleminin NP-zor (NP-hard) bir 
problem oldugu vurgulanmis ve farkli yontemler degerlendirilmistir:

\begin{itemize}
    \item \textbf{Tam Sayili Programlama:} Kucuk olcekli problemler 
    icin optimal cozumler uretebilen yontem.
    \item \textbf{Sezgisel Algoritmalar:} Buyuk olcekli problemlerde 
    kabul edilebilir surede iyi cozumler ureten yaklasimlar.
    \item \textbf{Metasezgisel Yontemler:} Genetik algoritmalar, 
    tavlama benzetimi ve tabu arama gibi optimizasyon teknikleri.
\end{itemize}

\subsection{Mobil Tabanli Tesis Yonetimi}
\label{subsec:mobil-tesis-yonetimi}

Kim ve Park (2016), Mobile-based Facility Management System for 
Enhanced Building Maintenance calismalarinda, mobil cihazlar 
araciligiyla tesis bakim ve yonetim sureclerinin 
iyilestirilmesini arastirmislardir.

Arastirmanin temel bulgulari sunlardir:
\begin{enumerate}
    \item Mobil tabanli sistemler, bakim talebi olusturma suresini 
    ortalama \%65 oraninda azaltmaktadir.
    \item Fotograf ve konum bilgisi ekleme ozelligi, sorunlarin 
    daha hizli ve dogru tespit edilmesini saglamaktadir.
    \item Push bildirim sistemi, acil durumlara yanit suresini 
    onemli olcude kisaltmaktadir.
\end{enumerate}

\subsection{Bulut Tabanli CMMS Sistemleri}
\label{subsec:bulut-cmms}

Wang ve Liu (2018), Cloud-Based Computerized Maintenance Management 
System (CMMS) for Small and Medium Enterprises calismalarinda, bulut 
tabanli bakim yonetim sistemlerinin tasarimini incelemislerdir.

\subsection{Personel Performans Degerlendirme Sistemleri}
\label{subsec:performans-degerlendirme}

Turkoglu ve Yilmaz (2020), Web Tabanli Personel Performans 
Degerlendirme Sistemi Tasarimi ve Uygulamasi baslikli calismalarinda, 
Turkiye'deki kamu kurumlari icin personel performans degerlendirme 
sistemi gelistirmislerdir.

Calismanin onemli katkilari sunlardir:
\begin{enumerate}
    \item Rol tabanli erisim kontrolu (RBAC) ile farkli kullanici 
    gruplarinin yetkilendirilmesi.
    \item Turk kamu kurumlarinin yapisina uygun hiyerarsik 
    organizasyon modeli.
    \item PHP ve MySQL tabanli web uygulamasi gelistirme deneyimi.
\end{enumerate}

\subsection{Hastane Temizlik Yonetim Sistemleri}
\label{subsec:hastane-temizlik}

Anderson ve Davies (2017), Digital Transformation in Hospital Cleaning 
Services: A Case Study isimli calismalarinda, bir buyuk hastanede 
temizlik hizmetlerinin dijitallestirilmesi surecini incelemislerdir.

Calismada, dijitallesme oncesi ve sonrasi performans metrikleri 
karsilastirilmistir:

\begin{table}[H]
\centering
\caption{Hastane Temizlik Dijitallesme Sonuclari}
\label{tab:hastane-dijitallesme}
\begin{tabular}{|l|c|c|c|}
\hline
\textbf{Metrik} & \textbf{Oncesi} & \textbf{Sonrasi} & \textbf{Iyilesme} \\
\hline
Gorev tamamlama orani & \%78 & \%95 & +17\% \\
\hline
Ortalama yanit suresi & 45 dk & 15 dk & -67\% \\
\hline
Musteri memnuniyeti & 3.2/5 & 4.5/5 & +40\% \\
\hline
\end{tabular}
\end{table}

\subsection{Universite Kampus Tesis Yonetimi}
\label{subsec:kampus-tesis}

Johnson ve Smith (2019), Integrated Facility Management System for 
University Campus calismalarinda, buyuk bir universite kampusunun 
tesis yonetimi icin butunlesik bir sistem gelistirmislerdir.

\subsection{Gorev Atama Algoritmalari}
\label{subsec:gorev-atama}

Ozkan ve Demir (2021), Genetik Algoritma Tabanli Dinamik Gorev 
Atama Sistemi calismalarinda, degisken kosullar altinda personele 
gorev atama probleminin cozumu icin genetik algoritma yaklasimi 
onermislerdir.

\subsection{PWA Tabanli Kurumsal Uygulamalar}
\label{subsec:pwa-kurumsal}

Nicholson ve arkadaslari (2019), Progressive Web Applications in 
Enterprise: Adoption and Benefits calismalarinda, PWA teknolojisinin 
kurumsal uygulamalarda kullanimini arastirmislardir.

Calismanin temel bulgulari sunlardir:
\begin{enumerate}
    \item PWAlar, yerel uygulamalara kiyasla \%60 daha dusuk 
    gelistirme maliyeti sunmaktadir.
    \item Tek kod tabani ile birden fazla platforma dagitim 
    mumkundur.
    \item Cevrimdisi calisma kapasitesi, saha calisanlari icin 
    kritik oneme sahiptir.
\end{enumerate}

\subsection{Rol Tabanli Erisim Kontrolu Sistemleri}
\label{subsec:rbac-sistemleri}

Sandhu ve arkadaslari (1996), temel RBAC (Role-Based Access Control) 
modelini tanimlayan klasik calismalarinda, rol tabanli erisim 
kontrolunun teorik temellerini olusturmuslardir.

\section{Ticari Yazilim Cozumleri}
\label{sec:ticari-yazilimlar}

Piyasada temizlik ve tesis yonetimi icin cesitli ticari yazilim 
cozumleri bulunmaktadir.

\begin{table}[H]
\centering
\caption{Uluslararasi Temizlik Yonetim Yazilimlari Karsilastirmasi}
\label{tab:uluslararasi-yazilimlar}
\begin{tabular}{|p{2.5cm}|p{3cm}|p{3cm}|p{3cm}|}
\hline
\textbf{Yazilim} & \textbf{Ozellikler} & \textbf{Avantajlar} & \textbf{Dezavantajlar} \\
\hline
CleanTelligent & Gorev yonetimi, Mobil uygulama & Kapsamli, Kullanici dostu & Yuksek maliyet \\
\hline
Swept & Cizelgeleme, Zaman takibi & Basit arayuz & Sinirli ozellestirme \\
\hline
Janitorial Manager & Musteri yonetimi, Faturalama & Entegre cozum & Karmasik kurulum \\
\hline
\end{tabular}
\end{table}

\section{KTTYS'nin Farklilasmasi}
\label{sec:farkliliklar}

Gelistirilen KTTYS sistemi, mevcut cozumlerden asagidaki 
noktalarda farklilasmaktadir:

\begin{enumerate}
    \item \textbf{Turkce Arayuz:} Tum sistem bilesenleri 
    Turkce olarak tasarlanmistir.
    \item \textbf{Turk Kamu Kurumu Yapisina Uygunluk:} 
    Fakulte bazli organizasyon ve Turk resmi tatil entegrasyonu.
    \item \textbf{Dusuk Maliyet:} Acik kaynak teknolojiler 
    kullanilarak lisans maliyeti ortadan kaldirilmistir.
    \item \textbf{Kolay Kurulum:} SQLite tabanli mimari ile 
    karmasik veritabani kurulumu gerektirmez.
    \item \textbf{Ozellestirilebilirlik:} Acik kaynak kod ile 
    kuruma ozel gelistirmeler yapilabilir.
    \item \textbf{PWA Teknolojisi:} Hem masaustu hem mobil 
    cihazlardan erisim tek uygulama ile saglanir.
\end{enumerate}
